%\section{Introduction}
Charged-current neutrino quasi-elastic scattering,
$\nu_\mu n\to\mu^-p$, distinguishes neutrino flavor and is valuable for 
neutrino oscillation experiments at energies near $1$~GeV
where it is responsible for a large fraction of the total reaction
cross-section~\cite{AguilarArevalo:2010wv,AguilarArevalo:2008rc,Abe:2011sj,Ayres:2004js}. For free nucleons the scattering process may be described by the standard theory of weak interactions with the inclusion of nucleon form factors~\cite{LlewellynSmith:1971zm}.   
Electron scattering~\cite{Bradford:2006yz} and neutrino scattering on deuterium~\cite{Bodek:2007vi,Kuzmin:2007kr} determine the most important form factors with good precision~\cite{Day:2012gb}.
However, neutrino oscillation experiments typically use detectors made of
heavier nuclei such as carbon~\cite{AguilarArevalo:2008qa,Ayres:2004js}, 
oxygen~\cite{Abe:2011ks}, iron~\cite{Michael:2008bc},  
or argon~\cite{Chen:2007ae, LBNE} where interactions with nucleons are 
modified by the nuclear environment.
These effects are commonly modeled using a 
relativistic Fermi gas~\cite{Smith:1972xh,Bodek:1980ar} 
(RFG) description of the nucleus as quasi-free, independent
nucleons with Fermi motion in a uniform binding potential.
Neutrino interaction generators~\cite{Andreopoulos201087,Hayato:2009zz,Golan:2012wx,Buss:2011mx,Lalakulich:2012gm} additionally simulate interactions of final state hadrons inside the target nucleus. The MiniBooNE experiment recently observed that this 
prescription, utilizing the free deuterium value for the axial form factor, does not accurately describe their
measurements of quasi-elastic scattering of neutrinos and
antineutrinos on a hydrocarbon target~\cite{AguilarArevalo:2007ab,AguilarArevalo:2013hm}.

The RFG approach may be supplemented by accounting for correlations between 
nucleons within the nucleus. Evidence for these correlations 
has been observed in electron-nucleus scattering~\cite{Subedi:2008zz}.
%Calculations using the random phase approximation 
%(RPA)~\cite{Graczyk:2003ru} predict a suppression of the quasi-elastic 
%cross-section at low four-momentum transfer squared ($Q^2$).  
Processes that produce multiple final state nucleons are thought to lead to enhancements in the cross-section~\cite{Carlson:2001mp,Shen:2012xz,Bodek:2011ps}.
These contributions are modeled using
different approaches~\cite{Martini:2010ex,Martini:2013sha,Nieves:2005rq} which
produce qualitatively similar though not quantitatively identical results.
The RFG model may also be replaced by an alternate spectral 
function (SF) model that calculates the joint probability distribution of 
scattering off a nucleon of given momentum and binding energy 
inside a nucleus~\cite{Benhar:1994hw}.
These nuclear effects may be significant for oscillation experiments 
seeking to measure the neutrino mass hierarchy 
and CP violation~\cite{Nieves:2012yz,Lalakulich:2012hs,Martini:2012uc}.

In this Letter we report the first study of muon neutrino
quasi-elastic interactions at energies between $1.5$ and $10$~GeV from the \minerva experiment, which uses a
finely segmented scintillator detector at Fermilab to measure muon
neutrino interactions on 
nuclear targets.  The analysis technique is similar to the one
employed to study the antineutrino reaction~\cite{nubarprl}.  The signal
has a $\mu^-$ in the final state along with one or more nucleons
(typically with a leading proton), and no mesons.  We reject events in
which mesons are produced by requiring that the hadronic system
recoiling against the muon has a low energy. That energy is measured
in two spatial regions. The {\it vertex energy} region corresponds to
a sphere around the vertex with a radius sufficient to contain a
proton (pion) with 225 (100)~MeV kinetic energy. This region is
sensitive to low energy protons which could arise from correlations
among nucleons in the initial state or final state interactions of the outgoing
hadron inside the target nucleus~\cite{Sobczyk:2012ms}. 
We do not use the vertex energy in
the event selection but study it for evidence of multi-nucleon
processes.  The {\it recoil energy} region includes energy depositions
outside of the vertex region and is sensitive to pions and higher
energy nucleons. We use the recoil energy to estimate and remove
inelastic backgrounds.

